% --- KOPIEER DIT BESTAND VOOR ELKE NIEUWE DAG ---

\subsection*{Wanneer \& wie?: [03-06-2026] - [Jonas]}

\textbf{Wat heb je gedaan?}\\
Gestart met het coderen voor de data reduction, ik maak gebruik van https://www.astropy.org/ccd-reduction-and-photometry-guide/v/dev/notebooks/00-00-Preface.html, waarbij we eerst de masterflats, -bias en -darks maken. Daarna calibreren we de foto's van SZ LYN om de benodigde data te extraheren.
\vspace{0.3cm} 

\noindent \textbf{Waarom heb je dat gedaan?}\\
Wij hebben voor ons onderzoek enkel de magnitude/licht nodig aangezien wij het LTTE meten, hierbij maken wij een O-C diagram waarbij we dus enkel licht meten. Hierdoor hebben we geen moeilijke code nodig, we moeten alleen reducen en hier de magnitude uithalen.

\noindent \textbf{Dingen die noemenswaardig zijn:}\\
Ik ben ook begonnen met het installeren van Marin. En wij hebben vanochtend even gecommuniceerd wat het plan is voor vandaag/morgen, gezien ik morgen om 18u een hertentamen heb.


\vspace{0.5cm}
\hrule 
\vspace{0.5cm}


